\chapter{Where this leaves us: an architecture recommendation}
\label{ch:recommendation}

This chapter turns the findings into a recommendation for the network to carry into the
empirical manuscript and to standardise on for the cross-task generalisation goal. We state
the decision criteria, score the candidates against them, and give a contingent
recommendation that names what must be resolved before it is final.

\section{Decision criteria}
A network is a good choice for this program to the extent that it:
\begin{enumerate}
\item \textbf{reproduces the behavioural signatures} --- the reliability-scaled cueing
benefit and the causal attention-to-sensitivity link (Chapter~\ref{ch:deepdive});
\item \textbf{exposes an interpretable, causally manipulable attention signal} --- the
attention must be readable and a genuine lever on behaviour (the microstimulation-like
causal result, Chapter~\ref{ch:deepdive}). Note the matched family showed that on the
short task the cueing itself is carried by \emph{memory}, not expressed as a cue-orienting
map (Chapter~\ref{ch:attention}), so this criterion is about a manipulable attention
\emph{lever}, with cue-orienting treated as a regime-dependent property, not a requirement;
\item \textbf{generalises across the task battery} (validity, set-size, value, distractor,
reward-structure) with the same architecture and harness;
\item \textbf{is neuro-plausible} --- recurrent gain modulation, a capacity-limited working
memory, reward-driven learning, and a manipulable priority signal; and
\item \textbf{trains reliably} under the stable harness with the frozen-VAE front-end.
\end{enumerate}

\section{Two non-negotiables}
Two components are settled by the evidence and should be fixed regardless of the remaining
choices. First, the \textbf{frozen pretrained VAE front-end} (Chapter~\ref{ch:perception}):
perception cannot be learned from reward by supervised reconstruction on the short task, the
VAE is task-parameter-invariant, and freezing it makes architecture comparisons clean.
Chapter~\ref{ch:jepa} shows the requirement is more precisely \emph{a front-end that reliably
encodes the task variable}, satisfiable by a frozen pretrained VAE \emph{or} by a
self-supervised V-JEPA predictive auxiliary (with anti-collapse machinery) paired with
frame-repeat or a convolutional encoder; the frozen VAE remains the recommended default for
the empirical manuscript on grounds of simplicity and characterisation. Second, the \textbf{stable
harness} --- PAC actor, distributional critic, prioritised episode replay, a time-lagged
target, and burn-in --- which is what makes a bootstrapped recurrent actor--critic learn on
a task with a $0.5$ harbor.

\section{The attention-mechanism choice (now decided)}
The matched family (Chapter~\ref{ch:attention}, Table~\ref{tab:familyresults}) was built to
settle this, and it did --- by eliminating the premise that one wiring would restore
cue-orienting attention. Reading the family against the criteria:
\begin{itemize}
\item The \textbf{pure multiplicative single-stream} form is the most faithful to the
published equation and reproduces the behaviour (correct $0.804$, hit $0.705$); its
attention is change-locked (spike $+0.143$, orient $\approx0$). On the new criterion (2) it
passes --- the attention is a readable, causally manipulable change detector --- it simply
does not orient to the cue, which the family shows \emph{nothing} on this task does.
\item The \textbf{two-LSTM} variant is the same paper-faithful multiplicative attention with
a second stacked memory feeding the heads, and it is the \emph{best detector in the family}
(correct $0.872$, hit $0.836$), while keeping the same interpretable, change-locked,
causally-manipulable attention. It is the strongest candidate on criteria (1), (4), (5) and
the new (2), at a minimal, well-motivated departure from the published network.
\item The \textbf{cross-attention} and \textbf{dual-memory} variants work (correct $0.840$,
$0.828$) but attend \emph{below} uniform at the cued image location (cue-orient $-0.179$,
$-0.174$): they route processing through memory keys rather than the image, which neither
helps behaviour over two-LSTM nor yields the cue-orienting map the lineage showed on the
older long-task setup. The earlier expectation that cross-attention would orient here was
not borne out on the controlled frozen-VAE short-task base.
\item The \textbf{affine} modulation \emph{failed} (always-wait collapse, hit $0$); it is
out. The \textbf{dual-head} sensory head is the only positive cue-orient ($+0.038$) but far
too weak to feature.
\item The \textbf{split priority/value (cross-talk)} architecture remains the vehicle for
mechanistic dissociations (Chapter~\ref{ch:crosstalk}); its $\mu/q$ mapping onto anatomy is
delicate, and the empirical paper's scope (per James) excludes architecture variants. It is
the subject of a \emph{separate} paper, not the recommended manuscript network.
\end{itemize}

\section{Recommendation}
\textbf{For the empirical manuscript:} a single-stream recurrent vision transformer with
(i) the frozen pretrained VAE front-end, (ii) the \textbf{multiplicative recurrent attention
block with a second stacked xLSTM memory} (the \code{two-lstm} variant) --- the best detector
in the matched family and a minimal, paper-faithful change to the published equation ---
(iii) a spatial xLSTM working memory, and (iv) a distributional actor--critic readout,
trained by the stable harness across the full task battery. This satisfies criteria
(1), (3), (4), (5) and the reframed (2): the attention is interpretable and a causal lever
on sensitivity, while the cueing benefit is carried, honestly and demonstrably, by working
memory.

The accompanying scientific claim is that, on the seven-step task, \textbf{the cueing
benefit is memory-borne and the attention map is a bottom-up change detector} --- a result
now established across six wirings, not assumed. The manuscript's attention figure should
therefore be the \emph{causal} one (biasing attention shifts sensitivity, with no induced
false alarms; Chapter~\ref{ch:deepdive}), not a cue-orienting map the mechanism does not
produce. Whether a longer delay or harder cue \emph{forces} pre-orienting --- as the lab's
older long-task networks displayed --- is a regime question to test with this same network,
not a reason to change the architecture.

\textbf{For the separate architecture line:} the split priority/value cross-talk network is
the vehicle for the coupling-required and double-dissociation results --- a strong paper in
its own right, kept out of the empirical manuscript by design.

\section{What must be resolved before this is final}
The attention-mechanism question is now settled; what remains: (1) the chosen network
(\code{two-lstm}) must be shown to generalise across at least the value (\code{vda4}) and
set-size (\code{setsize9}) tasks with no architecture change; (2) the cross-talk coupling
result must be replicated on the short task with the frozen VAE to confirm it survives the
perception fix; and (3) the regime question --- does a longer delay restore cue-orienting
attention in the recommended network? --- should be tested, since it bears on whether the
published orienting figure is reproducible at all in our setup or is specific to the long
task. Frozen VAE, the stable harness, and now the two-LSTM multiplicative attention are
fixed; the open items are generalisation and two targeted controls.
